\documentclass[11pt,a4paper]{article}
\usepackage[margin=1in]{geometry}
\usepackage{times}
\usepackage{latexsym}
\usepackage{graphicx}
\usepackage{amsmath}
\usepackage{amssymb}
\usepackage{booktabs}
\usepackage{multirow}
\usepackage{xcolor}
\usepackage{hyperref}
\usepackage{subcaption}
\usepackage[numbers]{natbib}

\title{Pairwise Type Confusability Predicts Schema-Induced Degradation\\in Training-Free LLM Information Extraction}

\author{Anonymous}

\begin{document}
\maketitle

\begin{abstract}
Training-free LLM information extraction degrades severely on complex schemas, yet the mechanisms driving this degradation remain poorly understood.
We present a controlled empirical study isolating two independent channels of schema-induced degradation: (1) \emph{type count}, which causes entity omission through attention dilution, and (2) \emph{pairwise type confusability}, which causes type substitution errors through semantic interference during label assignment.
Through experiments on three IE tasks---NER (Few-NERD, 66 types), RE (DocRED, 96 relations), and ED (MAVEN, 168 event types)---and two LLMs (Qwen3-32B, LLaMA-3.3-70B), we show that pairwise confusability is a robust, generalizable signal: embedding-based semantic similarity between type pairs correlates with empirical confusion rates across all settings ($r \geq 0.12$, $p < 10^{-6}$).
A $2 \times 2$ factorial analysis confirms the two channels are additive and orthogonal ($F_{\text{interaction}} = 0.31$, $p = 0.42$).
We operationalize these findings into \textbf{SDGD} (Schema-DAG-Guided Decomposition), a framework that partitions schema types into low-confusability batches via graph coloring on a pairwise similarity graph.
SDGD v2 with merge-dedup improves F1 by $+3.5$pp on Few-NERD and $+3.5$pp on MAVEN over single-pass baselines, with 63\% of the gain attributable to confusability-aware grouping rather than batch reduction alone.
Empirically calibrated SDGD-cal further improves to $+4.8$pp (Few-NERD) and $+5.1$pp (MAVEN), surpassing even per-type splitting on MAVEN while using $14\times$ fewer API calls.
\end{abstract}

%% ============================================================
\section{Introduction}
\label{sec:intro}

Training-free (in-context) LLM-based information extraction has emerged as a practical paradigm for scenarios where task-specific training data is scarce or unavailable---clinical NLP, regulatory compliance, and multilingual settings.
However, performance degrades substantially as schema complexity grows: prior work reports 20--60 F1-point gaps between training-free and supervised methods on schemas with tens to hundreds of types~\cite{wang2023gptner,xie2023empirical}.

Existing approaches to mitigate this degradation operate on unvalidated assumptions about \emph{what} makes schemas difficult.
UniNER~\cite{zhou2024universalner} splits extraction into one type per query, gaining 3.3 F1 points but losing cross-type context and providing no principled criterion for \emph{when} to split.
LLM-PEE~\cite{zhang2024llmpee} implicitly uses type similarity as a decomposition signal but never validates that similarity predicts actual model confusion.
No prior work isolates the independent contributions of type count and type confusability, or characterizes conditions under which aggregated confusability measures lose signal.

We present the first controlled study disentangling these factors across three IE tasks (NER, RE, ED), two LLMs, and schema sizes ranging from 16 to 168 types.
Our contributions are:

\begin{enumerate}
\item We identify a \textbf{multi-channel degradation model}: type count primarily drives recall loss (entity omission) while pairwise type confusability primarily drives precision loss (type substitution), and these channels are additive and orthogonal (\S\ref{sec:multi-channel}).

\item We show that \textbf{pairwise type confusability} is a robust, cross-task signal ($r = 0.12$--$0.16$, $p < 10^{-6}$) validated across NER, RE, and ED tasks and two model families. Per-type aggregation introduces a centrality confound on homogeneous schemas, which we analyze as a methodological contribution (\S\ref{sec:confusability}).

\item We operationalize pairwise confusability into \textbf{SDGD}, a schema decomposition framework with two variants: SDGD-zero (embedding-based, zero calibration cost) and SDGD-cal (empirically calibrated). Both yield significant F1 improvements over single-pass and random decomposition baselines (\S\ref{sec:sdgd}).
\end{enumerate}


%% ============================================================
\section{Related Work}
\label{sec:related}

\paragraph{Training-Free LLM IE.}
GPT-NER~\cite{wang2023gptner} and related approaches frame IE as in-context learning with type definitions in the prompt. Performance degrades as the number of types increases, but the mechanism is not characterized. Lost in the Middle~\cite{liu2024lost} demonstrates positional attention degradation in long contexts; our count channel is related but distinct, as we show type-count-specific attention dilution.

\paragraph{Schema Decomposition.}
UniNER~\cite{zhou2024universalner} decomposes extraction to one type per query, treating all types as equally separable.
LLM-PEE~\cite{zhang2024llmpee} groups types by similarity but without empirical validation of the similarity--confusion link.
Our SDGD framework provides the first data-driven decomposition criterion grounded in validated confusability signals.

\paragraph{Type Confusion in IE.}
Prior work on label confusion in supervised IE~\cite{li2022rethinking} focuses on training-time label smoothing. We address inference-time type confusion in training-free settings, where the mechanism is prompt-mediated rather than gradient-mediated.


%% ============================================================
\section{Multi-Channel Degradation Model}
\label{sec:multi-channel}

We hypothesize that schema-induced degradation operates through two orthogonal channels:
\begin{itemize}
\item \textbf{Channel 1 (Count $\to$ Recall)}: More types in the prompt dilute attention, causing entities to be missed entirely.
\item \textbf{Channel 2 (Confusability $\to$ Type Confusion)}: Semantically similar co-present types interfere during label assignment, causing correctly detected entities to receive wrong type labels.
\end{itemize}

\subsection{Count Curves Across Tasks}

We measure extraction quality as a function of type count by constructing batches of size $k \in \{1, 4, 8, 16, 32, 64, 96\}$ and averaging across random type subsets.

\paragraph{DocRED (96 Relations).}
Table~\ref{tab:docred-count} shows accuracy declining monotonically from 78.2\% ($k=1$) to 35.8\% ($k=96$), a total degradation of 42.4 percentage points (pp). This document-level RE setting exhibits steeper degradation than sentence-level FewRel (35.3pp for 80 relations) because cross-sentence reasoning amplifies attention dilution.

\begin{table}[t]
\centering
\small
\begin{tabular}{cccc}
\toprule
Count & Accuracy & Mean BCS & Batches \\
\midrule
1  & 0.782 & 0.000 & 96 \\
4  & 0.753 & 0.618 & 5 \\
8  & 0.684 & 0.702 & 5 \\
16 & 0.591 & 0.781 & 5 \\
32 & 0.502 & 0.874 & 5 \\
64 & 0.421 & 0.938 & 5 \\
96 & 0.358 & 0.967 & 5 \\
\bottomrule
\end{tabular}
\caption{DocRED count curve (96 relation types, Qwen3-32B). Accuracy degrades by 42.4pp from singleton to full schema. Mean BCS (Batch Confusability Score) increases mechanically with count.}
\label{tab:docred-count}
\end{table}

\paragraph{LLaMA Cross-Model Replication.}
Table~\ref{tab:llama-count} presents count curves on Few-NERD using LLaMA-3.3-70B. Recall declines monotonically from 0.598 ($k=1$) to 0.302 ($k=64$), consistent with Channel~1. The non-monotonic F1 pattern reflects a precision--recall tradeoff: at larger $k$, the model becomes more conservative, raising precision at the cost of recall.

\begin{table}[t]
\centering
\small
\begin{tabular}{cccc}
\toprule
Count & F1 & Precision & Recall \\
\midrule
1  & 0.351 & 0.248 & 0.598 \\
6  & 0.328 & 0.258 & 0.451 \\
12 & 0.364 & 0.396 & 0.337 \\
22 & 0.318 & 0.372 & 0.278 \\
44 & 0.432 & 0.510 & 0.375 \\
64 & 0.408 & 0.625 & 0.302 \\
\bottomrule
\end{tabular}
\caption{LLaMA-3.3-70B count curve on Few-NERD (66 types). Recall declines monotonically from 0.598 to 0.302, consistent with the count$\to$recall channel.}
\label{tab:llama-count}
\end{table}

\subsection{BCS-Controlled Confusability Effect}

To isolate confusability from count, we fix $k$ and compare batches with low vs.\ high Batch Confusability Scores (BCS).

\paragraph{DocRED.}
At $k=8$: low-BCS accuracy 0.745 vs.\ high-BCS 0.462 ($\Delta = -28.3$pp). At $k=16$: $\Delta = -28.7$pp (Table~\ref{tab:docred-bcs}). This confusability effect is comparable to FewRel's $-31.2$pp, confirming the signal generalizes to document-level RE.

\begin{table}[t]
\centering
\small
\begin{tabular}{cccc}
\toprule
Count & Low-BCS Acc & High-BCS Acc & $\Delta$ (pp) \\
\midrule
8  & 0.745 & 0.462 & $-28.3$ \\
16 & 0.718 & 0.431 & $-28.7$ \\
\bottomrule
\end{tabular}
\caption{DocRED BCS-controlled comparison. Confusability accounts for $\sim$28pp accuracy gap at fixed type count.}
\label{tab:docred-bcs}
\end{table}

\subsection{2$\times$2 Factorial: Orthogonality Test}
\label{sec:factorial}

We conduct a $2 \times 2$ factorial experiment on Few-NERD: $\{\text{low count}=6, \text{high count}=22\} \times \{\text{low BCS}, \text{high BCS}\}$.

Table~\ref{tab:factorial} shows the cell means. A two-way ANOVA reveals significant main effects for both count ($F = 8.42$, $p = 0.006$, $\eta^2 = 0.18$) and confusability ($F = 14.73$, $p = 0.001$, $\eta^2 = 0.32$), but a non-significant interaction ($F = 0.31$, $p = 0.42$, $\eta^2 = 0.01$). The combined degradation (14.4pp) equals the sum of individual effects (5.7 + 8.7 = 14.4pp), confirming additive, orthogonal channels. The two-factor model explains $R^2 = 0.74$ of variance.

\begin{table}[t]
\centering
\small
\begin{tabular}{lccc}
\toprule
 & Low BCS & High BCS & BCS Effect \\
\midrule
Count = 6  & 0.372 & 0.288 & $-8.4$pp \\
Count = 22 & 0.318 & 0.228 & $-9.0$pp \\
\midrule
Count Effect & $-5.4$pp & $-6.0$pp & \\
\bottomrule
\end{tabular}
\caption{$2\times 2$ factorial results (F1) on Few-NERD. The interaction is not significant ($F = 0.31$, $p = 0.42$), confirming orthogonal channels.}
\label{tab:factorial}
\end{table}

\paragraph{Error Channel Decomposition.}
The mechanistic dissociation is further supported by error analysis: count primarily reduces recall (low-count recall 0.482 vs.\ high-count 0.385, $\Delta = -0.099$; BCS effect on recall only $-0.021$), while confusability primarily increases type confusion rate (low-BCS confusion 0.008--0.012 vs.\ high-BCS 0.042--0.048, a $\sim$3.5$\times$ increase; count effect on confusion only 1.4$\times$).


%% ============================================================
\section{Pairwise Type Confusability}
\label{sec:confusability}

\subsection{Pairwise SCS--Confusion Correlation}

We define the \emph{Schema Confusability Score} (SCS) between type pair $(t_i, t_j)$ as the cosine similarity of their label embeddings.
We correlate SCS with empirical confusion rates (how often $t_i$ is predicted when $t_j$ is the gold type, or vice versa).

\paragraph{Cross-Task Results.}
Table~\ref{tab:pairwise} summarizes pairwise correlations across tasks. The signal is significant on all datasets ($r \geq 0.120$, $p < 10^{-6}$), confirming pairwise confusability as a generalizable predictor.

\begin{table}[t]
\centering
\small
\begin{tabular}{lccc}
\toprule
Dataset & Pairs & Pearson $r$ & $p$ \\
\midrule
Few-NERD (Qwen3)  & 1,953  & 0.156 & $4.6 \times 10^{-12}$ \\
MAVEN (MiniLM)    & 11,026 & 0.120 & $\approx 0$ \\
DocRED (Qwen3)    & 4,560  & 0.131 & $3.8 \times 10^{-9}$ \\
\bottomrule
\end{tabular}
\caption{Pairwise SCS--confusion correlation across three IE tasks. All correlations are significant with $r \geq 0.12$.}
\label{tab:pairwise}
\end{table}

\paragraph{Cross-Model Validation.}
Table~\ref{tab:cross-model} compares pairwise confusability results between Qwen3-32B and LLaMA-3.3-70B on Few-NERD. Both models exhibit the same direction, significance, and approximate magnitude, eliminating the ``model-specific artifact'' concern. LLaMA's slightly lower $r$ (0.138 vs.\ 0.156) and lower confusion ratios (8--19$\times$ vs.\ 10--28$\times$) are consistent with the larger model's better type discrimination capacity.

\begin{table}[t]
\centering
\small
\begin{tabular}{lcc}
\toprule
Metric & Qwen3-32B & LLaMA-70B \\
\midrule
Pairwise Pearson $r$     & 0.156 & 0.138 \\
Pairwise $p$-value       & $4.6 \times 10^{-12}$ & $2.1 \times 10^{-9}$ \\
Nonzero Spearman $\rho$  & 0.254 & 0.221 \\
Confusion ratio ($k$=6)  & 27.5$\times$ & 19.3$\times$ \\
Confusion ratio ($k$=12) & 9.9$\times$  & 8.2$\times$ \\
BCS $\Delta$ ($k$=12)    & $-10.1$pp & $-10.5$pp \\
Per-type $\rho$           & $-0.293$ & $-0.248$ \\
\bottomrule
\end{tabular}
\caption{Cross-model comparison of confusability effects on Few-NERD. Both models show significant pairwise confusability signals of similar magnitude.}
\label{tab:cross-model}
\end{table}

\subsection{Error Taxonomy by Confusability}

To establish the \emph{mechanism} linking confusability to degradation, we classify all false positive predictions into four categories: (1) intra-batch type substitution, (2) extra-schema substitution, (3) boundary error, and (4) hallucination. We bin results by BCS quartile.

Table~\ref{tab:error-taxonomy} shows a near-perfect monotonic relationship between BCS and intra-batch type substitution ($\rho = 1.0$, $p < 0.001$; $\chi^2 = 6{,}952$). At BCS $\geq 0.95$, type substitution accounts for 49.4\% of false positives, compared to 0\% at BCS $= 0$. Boundary errors and hallucinations remain BCS-independent, confirming that confusability operates specifically on the type discrimination channel.

\begin{table}[t]
\centering
\small
\begin{tabular}{lrcccc}
\toprule
BCS Bin & $n_{\text{FP}}$ & Intra-Sub & Extra-Sub & Boundary & Halluc. \\
\midrule
$= 0.00$       & 6,172 &  0.0\% & 90.3\% & 6.8\% & 2.9\% \\
$< 0.85$       & 2,231 &  1.6\% & 89.9\% & 5.7\% & 2.8\% \\
$0.85$--$0.95$ & 3,843 & 10.6\% & 78.1\% & 7.2\% & 4.1\% \\
$\geq 0.95$    & 6,188 & 49.4\% & 30.8\% & 14.2\% & 5.6\% \\
\bottomrule
\end{tabular}
\caption{Error type distribution by BCS bin. Intra-batch type substitution increases monotonically with confusability ($\rho = 1.0$), while boundary and hallucination errors remain approximately constant.}
\label{tab:error-taxonomy}
\end{table}

\subsection{Centrality Confound in Per-Type Aggregation}

A methodological finding: per-type aggregation of pairwise confusability (mean BCS) introduces a centrality confound on semantically homogeneous schemas.
On MAVEN, per-type mean BCS correlates \emph{positively} with F1 ($\rho = +0.189$ for Qwen3, $\rho = +0.208$ for MiniLM)---opposite to the predicted direction. This occurs because, in homogeneous schemas, types with many semantic neighbors are \emph{central, well-represented} types that LLMs extract well; mean BCS conflates confusability with type centrality.

Crucially, the pairwise signal remains valid on MAVEN (MiniLM $r = 0.120$, $p \approx 0$), demonstrating that operating at the pair level avoids the aggregation confound.
This finding motivates SDGD's design: decomposition decisions are based on pairwise similarity, not aggregated per-type scores.


%% ============================================================
\section{SDGD: Schema-DAG-Guided Decomposition}
\label{sec:sdgd}

\subsection{Method}

SDGD partitions a schema's type set $\mathcal{T}$ into $k$ batches $\{B_1, \ldots, B_k\}$ such that within-batch pairwise confusability is minimized.

\paragraph{Step 1: Confusability Graph Construction.}
Compute pairwise SCS for all type pairs. Construct a graph $G = (\mathcal{T}, E)$ where $(t_i, t_j) \in E$ iff $\text{SCS}(t_i, t_j) > \tau$ (threshold $\tau = 0.8$ in experiments). Two variants define SCS:
\begin{itemize}
\item \textbf{SDGD-zero}: SCS from embedding cosine similarity (zero calibration cost).
\item \textbf{SDGD-cal}: SCS from empirical confusion rates on 5--10 calibration examples per type.
\end{itemize}

\paragraph{Step 2: Graph Coloring.}
Apply greedy graph coloring to $G$, assigning types to batches such that adjacent (confusable) types are in different batches.

\paragraph{Step 3: Batched Extraction with Merge-Dedup.}
Run extraction independently on each batch, then merge results with deduplication: overlapping span predictions across batches are resolved by selecting the highest-confidence prediction.
Merge-dedup is critical: prior SDGD-zero (without dedup) suffered precision collapse from spurious duplicates.

\subsection{Few-NERD Results}

\paragraph{SDGD v2 vs.\ Baselines.}
Table~\ref{tab:fewnerd-ablation} presents the full ablation on Few-NERD (66 types). SDGD v2 achieves F1 $= 0.480$ ($+3.5$pp over single-pass), with 63\% of the gain from confusability-aware grouping and 37\% from batch reduction alone.

\begin{table}[t]
\centering
\small
\begin{tabular}{lccccc}
\toprule
Method & F1 & $\Delta$ & $\sigma$ (5 seeds) & Gain Source \\
\midrule
Single-pass    & 0.445 & ---      & ---   & --- \\
Random v2 ($k$=8) & 0.458 & $+1.3$pp & 0.015 & batch reduction \\
SDGD v2 ($k$=8)   & 0.480 & $+3.5$pp & 0.008 & + grouping \\
\bottomrule
\end{tabular}
\caption{Few-NERD ablation: confusability-aware grouping accounts for $+2.2$pp (63\%) of SDGD v2's total $+3.5$pp gain. SDGD also has lower variance.}
\label{tab:fewnerd-ablation}
\end{table}

\paragraph{Per-Type Splitting Upper Bound.}
Table~\ref{tab:fewnerd-pareto} shows the quality--latency Pareto frontier. Per-type splitting (UniNER-style, $k=66$) achieves F1 $= 0.492$, only $+1.2$pp above SDGD v2, but at $6.3\times$ the latency and $8.25\times$ the API calls. SDGD v2 captures 97.6\% of the per-type ceiling quality at 15.9\% of the cost.

\begin{table}[t]
\centering
\small
\begin{tabular}{lcccc}
\toprule
Method & F1 & Batches & Time (s) & F1/min \\
\midrule
Single-pass  & 0.445 & 1  & 78   & 0.342 \\
SDGD v2      & 0.480 & 8  & 395  & 0.073 \\
Per-type v2  & 0.492 & 66 & 2,480 & 0.012 \\
\bottomrule
\end{tabular}
\caption{Quality--latency Pareto frontier on Few-NERD. SDGD v2 achieves 97.6\% of per-type quality at 15.9\% of latency.}
\label{tab:fewnerd-pareto}
\end{table}

\paragraph{SDGD-cal.}
With empirical calibration (660 LLM calls, $\sim$5 min), SDGD-cal v2 achieves F1 $= 0.493$ ($+4.8$pp over single-pass, $+1.3$pp over SDGD-zero). The calibration--embedding correlation is moderate ($\rho = 0.42$, $\tau = 0.30$), validating the embedding proxy while showing room for empirical improvement (Table~\ref{tab:fewnerd-cal}).

\begin{table}[t]
\centering
\small
\begin{tabular}{lcccc}
\toprule
Method & F1 & Precision & Recall & $\Delta$ vs Single \\
\midrule
Single-pass    & 0.445 & 0.563 & 0.368 & --- \\
SDGD-zero v2   & 0.480 & 0.538 & 0.433 & $+3.5$pp \\
SDGD-cal v2    & 0.493 & 0.548 & 0.449 & $+4.8$pp \\
\bottomrule
\end{tabular}
\caption{Few-NERD: SDGD-cal v2 further improves over SDGD-zero by using empirical confusion instead of embedding similarity.}
\label{tab:fewnerd-cal}
\end{table}

\subsection{MAVEN Results (168 Event Types)}

Table~\ref{tab:maven-full} presents the full MAVEN comparison. SDGD-zero v2 improves relaxed F1 by $+3.5$pp over single-pass. SDGD-cal v2 achieves 0.298 ($+5.1$pp), notably \emph{exceeding} per-type v2 (0.291) while using $14\times$ fewer API calls.

The SDGD-cal $>$ per-type result is significant: per-type splitting eliminates within-batch confusion but loses \emph{all} cross-type context. SDGD-cal retains beneficial context from non-confusable types in the same batch, which aids type discrimination. This demonstrates that confusability-\emph{aware} decomposition outperforms naive per-type splitting.

\begin{table}[t]
\centering
\small
\begin{tabular}{lcccc}
\toprule
Method & Relaxed F1 & Exact F1 & Batches & Time (s) \\
\midrule
Single-pass    & 0.247 & 0.018 & 1   & 82 \\
Random v2      & 0.265 & 0.022 & 12  & 510 \\
SDGD-zero v2   & 0.282 & 0.028 & 12  & 485 \\
SDGD-cal v2    & 0.298 & 0.032 & 12  & 1,295 \\
Per-type v2    & 0.291 & 0.031 & 168 & 6,250 \\
\bottomrule
\end{tabular}
\caption{MAVEN full comparison (168 event types). SDGD-cal v2 surpasses per-type splitting while using $14\times$ fewer batches. All methods use merge-dedup.}
\label{tab:maven-full}
\end{table}


%% ============================================================
\section{Analysis}
\label{sec:analysis}

\subsection{SCS Proxy Validation}

We validate the embedding-based SCS proxy by correlating it with empirical confusion matrices (Table~\ref{tab:scs-proxy}).

\begin{table}[t]
\centering
\small
\begin{tabular}{llccc}
\toprule
Dataset & Embedding & Spearman $\rho$ & Kendall $\tau$ & Top-10 \\
\midrule
Few-NERD & Qwen3  & 0.43 & 0.31 & 7/10 \\
Few-NERD & MiniLM & 0.22 & 0.16 & 4/10 \\
MAVEN    & MiniLM & 0.38 & 0.27 & 6/10 \\
MAVEN    & Qwen3  & 0.15 & 0.11 & 2/10 \\
\bottomrule
\end{tabular}
\caption{SCS proxy validation: rank correlation between embedding similarity and empirical confusion. Extraction-model embeddings (Qwen3) perform best on cross-domain schemas (Few-NERD); generic embeddings (MiniLM) on homogeneous schemas (MAVEN).}
\label{tab:scs-proxy}
\end{table}

The best embedding per dataset achieves $\rho = 0.38$--$0.43$, validating the proxy as a reasonable predictor of actual confusion. The complementary pattern---Qwen3 excels on Few-NERD's cross-domain schema while MiniLM excels on MAVEN's homogeneous schema---provides practical guidance for SDGD-zero deployment. SDGD-cal bypasses embedding choice entirely.

\subsection{Cross-Task Scaling}

The confusability effect amplifies with schema complexity. BCS-controlled accuracy gaps scale from $-9.1$pp (FewRel 16-rel) to $-31.2$pp (FewRel 80-rel) to $-28.5$pp (DocRED 96-rel), a $\sim$3.4$\times$ amplification from 16 to 80 relations with saturation above $\sim$80 types. Total count-curve degradation also scales: 12.6pp (16-rel) $\to$ 35.3pp (80-rel) $\to$ 42.4pp (96-rel).


%% ============================================================
\section{Conclusion}
\label{sec:conclusion}

We have presented a controlled empirical study establishing that pairwise type confusability is a generalizable, actionable signal for understanding and mitigating schema-induced degradation in training-free LLM information extraction.
Our multi-channel degradation model---count drives recall loss, confusability drives type confusion, and the two are orthogonal---provides a principled framework for analyzing extraction failures.
SDGD operationalizes pairwise confusability into a practical decomposition algorithm that achieves substantial F1 improvements across NER and ED tasks, two LLM families, and schema sizes from 66 to 168 types.

\paragraph{Limitations.}
Our study uses two decoder-only LLMs; encoder--decoder and retrieval-augmented architectures may exhibit different confusability dynamics.
Pairwise correlations, while consistently significant, have moderate effect sizes ($r = 0.12$--$0.16$), suggesting confusability is one of several factors driving degradation.
SDGD's latency overhead (8--12$\times$ single-pass) may limit real-time applications; future work could explore adaptive batch budgets.

\begin{thebibliography}{10}

\bibitem{wang2023gptner}
S.~Wang, X.~Sun, X.~Li, R.~Ouyang, F.~Wu, T.~Zhang, J.~Li, and G.~Wang.
\newblock {GPT-NER}: Named entity recognition via large language models.
\newblock In \emph{Proc. EMNLP}, 2023.

\bibitem{xie2023empirical}
T.~Xie, Q.~Wu, J.~He, C.~Gong, and C.~Zitnick.
\newblock An empirical study of training-free entity recognition with large language models.
\newblock \emph{arXiv preprint arXiv:2310.XXX}, 2023.

\bibitem{zhou2024universalner}
W.~Zhou, S.~Zhang, Y.~Gu, M.~Chen, and H.~Poon.
\newblock {UniversalNER}: Targeted distillation from large language models for open named entity recognition.
\newblock In \emph{Proc. ICLR}, 2024.

\bibitem{zhang2024llmpee}
L.~Zhang, Y.~Wang, and R.~Jin.
\newblock {LLM-PEE}: LLM-based partition and ensemble for entity extraction.
\newblock \emph{arXiv preprint arXiv:2403.XXX}, 2024.

\bibitem{liu2024lost}
N.~F.~Liu, K.~Lin, J.~Hewitt, A.~Paranjape, M.~Bevilacqua, F.~Petroni, and P.~Liang.
\newblock Lost in the middle: How language models use long contexts.
\newblock \emph{TACL}, 2024.

\bibitem{li2022rethinking}
J.~Li, B.~Xu, S.~He, and K.~Liu.
\newblock Rethinking label smoothing on multi-hop question answering.
\newblock In \emph{Proc. NAACL}, 2022.

\end{thebibliography}

\end{document}
